\documentclass[../../main.tex]{subfiles}
\begin{document}

\begin{flushright}
  \footnotesize\textit{【自動文字起こし・要確認】}
\end{flushright}

[解] (1) 選んだ頂点 $X_1$ とする．

\begin{figure}[htb]
  \centering
  \begin{tikzpicture}[scale=1.25, every node/.style={font=\small}]
    \coordinate (A) at (0,1.8);
    \coordinate (B) at (-1.55,-0.9);
    \coordinate (C) at (1.55,-0.9);
    \coordinate (O) at (0,0);
    \coordinate (Pzero) at (-0.40,-0.20);
    \coordinate (Pone) at ($(A)!0.5!(Pzero)$);
    \fill[blue!6] (A)--(B)--(C)--cycle;
    \draw[blue!75!black,very thick] (A)--(B)--(C)--cycle;
    \draw[teal!70!black,thick] (A)--(Pzero);
    \fill[gray!75!black] (O) circle (1pt);
    \fill[teal!70!black] (Pzero) circle (1.5pt);
    \fill[teal!70!black] (Pone) circle (1.5pt);
    \node[above] at (A) {$A=X_1$};
    \node[below left] at (B) {$B$};
    \node[below right] at (C) {$C$};
    \node[right] at (O) {$O$};
    \node[below left] at (Pzero) {$P_0$};
    \node[right] at (Pone) {$P_1$};
  \end{tikzpicture}
  \caption{選んだ頂点が$A$のとき，$P_1$は$AP_0$の中点}
\end{figure}

$\overrightarrow{OP_1} = \dfrac{1}{2}(\overrightarrow{OP_0} + \overrightarrow{OX_1})$

$|\overrightarrow{OP_1}|^2 = \dfrac{1}{4}\{|\overrightarrow{OP_0}|^2 + |\overrightarrow{OX_1}|^2 + 2\overrightarrow{OP_0} \cdot \overrightarrow{OX_1}\}$

$= \dfrac{1}{4}|\overrightarrow{OP_0}|^2 + 1 + 2\overrightarrow{OP_0} \cdot \overrightarrow{OX_1} \quad (\because |\overrightarrow{OX_1}| = 1)$

より，
\[ E_1 = \sum_{X_1=A,B,C} |\overrightarrow{OP_1}|^2 \cdot \dfrac{1}{3} \]
\[ = \dfrac{1+|\overrightarrow{OP_0}|^2}{4} + \dfrac{2}{3} \cdot \dfrac{1}{4} \overrightarrow{OP_0} \cdot (\overrightarrow{OA} + \overrightarrow{OB} + \overrightarrow{OC}) = \dfrac{1+|\overrightarrow{OP_0}|^2}{4} \quad (\because \overrightarrow{OA} + \overrightarrow{OB} + \overrightarrow{OC} = \vec{0}) \] \hfill(終)

\paragraph{(2)}

点$X$に対し，$\overrightarrow{OX} = \vec{x}$ とする．題意から，
\[ \overrightarrow{P_{n+1}} = \dfrac{1}{2}(\overrightarrow{P_n} + \overrightarrow{x_{n+1}}) \]
だから，くり返し用いて，
\[ \overrightarrow{P_n} = \dfrac{1}{2}\overrightarrow{x_n} + \dfrac{1}{4}\overrightarrow{x_{n-1}} + \cdots + \dfrac{1}{2^n}\overrightarrow{x_1} + \dfrac{1}{2^n}\overrightarrow{OP_0} \]
\[ = \dfrac{1}{2^n}\overrightarrow{OP_0} + \sum_{i=1}^n \dfrac{1}{2^{n+1-i}}\overrightarrow{OX_i} \] \hfill(終)

[解2]

\paragraph{(3)}

$E_1 = \dfrac{1}{4} \ (\because (1)) \cdots \text{③}$ である．ここで，$n \ge 2$ の時，
\[ \overrightarrow{P_n} = \dfrac{1}{2} \{ \overrightarrow{P_{n-1}} + \overrightarrow{x_n} \} \]
\[ |\overrightarrow{P_n}|^2 = \dfrac{1}{4}\{ |\overrightarrow{P_{n-1}}|^2 + |\overrightarrow{x_n}|^2 \} + \dfrac{1}{2}\overrightarrow{P_{n-1}} \cdot \overrightarrow{x_n} \cdots \text{④} \]
であり，
\[ E(|\overrightarrow{P_{n-1}}|) = E_{n-1} \]
\[ E(|\overrightarrow{x_n}|^2) = 1 \]
\[ E(\overrightarrow{P_{n-1}} \cdot \overrightarrow{x_n}) = \dfrac{1}{3}\overrightarrow{P_{n-1}} \cdot (\vec{a} + \vec{b} + \vec{c}) = 0 \]
を④に代入して
\[ E_n = \dfrac{1}{4} (E_{n-1} + 1) \]
③とあわせて
\[ E_n = \dfrac{1}{3}\{1 - \left(\dfrac{1}{4}\right)^n\} \] \hfill(終)

\paragraph{(3)}

$P_0 = O$ の時，
\[ \overrightarrow{OP_n} = \sum_{i=1}^n \dfrac{1}{2^{n+1-i}} \overrightarrow{OX_i} \]
だから，$Y_i$ を，$\overrightarrow{y_i} = \dfrac{1}{2^{n+1-i}} \overrightarrow{x_i}$ となるよう定めると，
\[ |\overrightarrow{P_n}|^2 = |\overrightarrow{y_1} + \overrightarrow{y_2} + \cdots + \overrightarrow{y_n}|^2 \]
\[ = \sum_{i=1}^n |\overrightarrow{y_i}|^2 + 2 \sum_{i < j} \overrightarrow{y_i} \cdot \overrightarrow{y_j} \]
\[ = \sum_{i=1}^n \dfrac{1}{4^{n+1-i}} + 2 \sum_{i < j} \overrightarrow{y_i} \cdot \overrightarrow{y_j} \quad (\because |\overrightarrow{x_i}| = 1) \]
\[ = \dfrac{1}{3}\{1 - (1/4)^n\} + 2 \sum_{i < j} \overrightarrow{y_i} \cdot \overrightarrow{y_j} \cdots \text{①} \]
である．ここで，$i \ne j$ に対して，
\[ \overrightarrow{x_i} \cdot \overrightarrow{x_j} =
  \begin{cases} 1 & (X_i = X_j) \\ -1/2 & (X_i \ne X_j)
\end{cases} \cdots \text{②} \]
だから $\overrightarrow{x_i} \cdot \overrightarrow{x_j}$ の期待値 $P_{i,j}$ は
\[ P = \dfrac{1}{3} \cdot 1 + \dfrac{2}{3} \left(-\dfrac{1}{2}\right) = 0 \]
なので，$\sum_{i < j} \overrightarrow{y_i} \cdot \overrightarrow{y_j}$ の期待値 $Q$ は，
\[ Q = \sum_{i < j} \dfrac{1}{2^{n+1-i} \cdot 2^{n+1-j}} \cdot P_{i,j} = 0 \]
だから ①より
\[ E_n = \dfrac{1}{3}(1 - (1/4)^n) + Q = \dfrac{1}{3} \{ 1 - \left(\dfrac{1}{4}\right)^n \} \] \hfill(終)

\end{document}
